\documentclass[12pt]{article}

% -------------------------------------------------
% Packages
% -------------------------------------------------
\usepackage[a4paper,margin=1in]{geometry}
\usepackage{graphicx}
\usepackage{float}
\usepackage{caption}
\usepackage{subcaption}
\usepackage{amsmath}
\usepackage{amssymb}
\usepackage{booktabs}
\usepackage{hyperref}
\usepackage{listings}
\usepackage{xcolor}

% -------------------------------------------------
% Listings setup (for SQL / code)
% -------------------------------------------------
% ==== SQL CODE SUPPORT ====
% Basic SQL listing style
\lstdefinestyle{sqlstyle}{
	language=SQL,
	basicstyle=\ttfamily\small,
	keywordstyle=\bfseries,
	commentstyle=\itshape,
	stringstyle=\ttfamily,
	columns=fullflexible,
	keepspaces=true,
	showstringspaces=false,
	frame=single,
	breaklines=true
}

% Environment for SQL blocks:
%   \begin{sqlcode}
	%   SELECT ...
	%   \end{sqlcode}
\lstnewenvironment{sqlcode}[1][]{
	\lstset{style=sqlstyle,#1}
}{}

% Inline SQL:
%   \sqlinline{SELECT * FROM Student;}
\newcommand{\sqlinline}[1]{\lstinline[language=SQL]!#1!}

% ============================================

% -------------------------------------------------
% Title information
% -------------------------------------------------
\title{
	\textbf{CAMM 535 – Fundamentals of Biological Databases} \\
	\vspace{0.5em}
	\textbf{Final Project – Report 2} \\
	\vspace{0.5em}
	Physical Design, Testing and Visualization
}

\author{
	Team Number: Four \\
	Disease: Schwannoma \\
	\vspace{0.5em} \\
	Student Name(s): Alireza Khoeini, Alireza Noroozi, Ghada Bahlul
}

\date{
	January 2026
}

\begin{document}
	
	\maketitle
	\newpage
	
	\section*{Part 1: Recap of Report 1 and Project Objectives}
	
	% =========================================================
	% Part 1 (Report 2): Recap of Report 1 (Objectives + Findings)
	% Paste this into the empty Part 1 section of Report 2.
	% (Assumes your preamble already loads graphicx, float, hyperref, etc.)
	% =========================================================
	
	\section{Part 1 --- Recap of Report 1: Objectives and Key Findings}
	
	\subsection{Overall objective of Report 1}
	The first report established the biological motivation and the full data-acquisition pipeline for our Schwannoma
	project, and translated that pipeline into a coherent conceptual database design.
	The main goal was to build a unified, phenotype-centered dataset that connects:
	(i) schwannoma-associated human genes, (ii) network-expanded genes based on protein--protein interactions,
	(iii) short genetic variants (SNPs/indels) mapped to those genes, (iv) transcript-level genomic coordinates and exon
	structure, (v) gene-to-protein mappings with functional protein attributes, and (vi) expression-level evidence from a
	relevant GEO dataset (GEO2R results).
	The outcome of Report 1 was therefore twofold:
	\begin{enumerate}
		\item a reproducible multi-database retrieval workflow (BioMart $\rightarrow$ STRING $\rightarrow$ BioMart variants $\rightarrow$ UCSC $\rightarrow$ UniProt $\rightarrow$ GEO/GEO2R),
		\item a schema-ready set of tables (and intermediate TSV files) that support later physical design, testing, and visualization tasks required in the final project.
	\end{enumerate}
	
	\subsection{Disease / phenotype context and motivation (Schwannoma)}
	Schwannomas are typically benign tumors that arise from Schwann cells, often along peripheral or cranial nerves.
	Even though they are not usually malignant, their location and compressive growth pattern can lead to significant
	clinical morbidity. A major and well-studied genetic driver is disruption of the \textit{NF2} tumor suppressor gene,
	which encodes merlin (schwannomin). Merlin participates in multiple growth-control pathways, including Hippo,
	PI3K/AKT, and MAPK signaling, so loss of NF2 function is a biologically meaningful entry point for studying
	schwannoma mechanisms and downstream pathway dysregulation.
	
	\begin{figure}[H]
		\centering
		\includegraphics[width=0.88\linewidth]{cbtrus_figure12.jpg}
		\caption{Background epidemiology and distribution context (CBTRUS, 2017--2021) used in Report 1 to motivate schwannoma as a prevalent non-malignant nerve sheath tumor entity.}
		\label{fig:cbtrus_context}
	\end{figure}
	
	\begin{figure}[H]
		\centering
		\includegraphics[width=0.88\linewidth]{nf2_merlin_pathway.png}
		\caption{Merlin (\textit{NF2}) signaling context summarized in Report 1. Merlin regulates multiple growth-control pathways; loss of merlin removes growth inhibition and promotes tumor formation.}
		\label{fig:nf2_merlin}
	\end{figure}
	
	\subsection{Key data acquisition pipeline and findings}
	
	\subsubsection{Step 1: Phenotype-driven gene identification using Ensembl BioMart}
	We began with a phenotype-first strategy rather than a manually curated list.
	Using Ensembl BioMart, we selected the \textbf{Human genes (GRCh38.p14)} dataset and applied the phenotype filter
	\textbf{schwannoma}. From the resulting set of genes, we extracted core gene-level attributes needed for downstream
	integration and database design, including:
	\textit{Ensembl Gene Stable ID}, gene name, chromosome/scaffold, and genomic coordinates (start/end).
	This step ensured that the initial gene list was grounded in Ensembl's phenotype annotation layer and provided a
	clean identifier backbone (Ensembl stable IDs) for later joins.
	
	\begin{figure}[H]
		\centering
		\includegraphics[width=0.92\linewidth]{biomart_gene_selection.png}
		\caption{BioMart configuration used in Report 1 for phenotype-based schwannoma gene retrieval and selection of gene-level attributes.}
		\label{fig:biomart_genes}
	\end{figure}
	
	\paragraph{Finding (Step 1).}
	The output of this step was a phenotype-associated schwannoma gene set in Ensembl identifiers, suitable as the
	``seed'' list for interaction-based expansion and variant retrieval.
	
	\subsubsection{Step 2: Interaction-based expansion using STRING}
	The BioMart-derived gene list was used as input to the STRING database to expand the set using protein--protein
	interaction (PPI) evidence. We used the \textbf{Multiple Proteins by Names / Identifiers} input mode and restricted
	the organism to \textbf{\textit{Homo sapiens}}. We constructed a \textbf{full STRING network} (functional + physical
	associations), enabled all evidence channels (text mining, experiments, curated databases, co-expression,
	neighborhood, gene fusion, co-occurrence), and set the minimum interaction score to \textbf{medium confidence (0.400)}.
	In accordance with the project requirement, we expanded using the \textbf{first shell} and capped the maximum number
	of interactors at \textbf{100}.
	
	\begin{figure}[H]
		\centering
		\includegraphics[width=0.92\linewidth]{string_gene_input.png}
		\caption{STRING multiple-protein input interface used in Report 1 (HGNC symbols).}
		\label{fig:string_input}
	\end{figure}
	
	\begin{figure}[H]
		\centering
		\includegraphics[width=0.92\linewidth]{string_hires_image.png}
		\caption{High-resolution STRING interaction network visualization used in Report 1.}
		\label{fig:string_network}
	\end{figure}
	
	\paragraph{Finding (Step 2).}
	The resulting STRING-expanded network contained \textbf{72 nodes} and showed dense connectivity among genes involved
	in pathways commonly implicated in tumor biology and schwannoma-associated signaling (including nodes consistent with
	growth control, receptor signaling, and downstream kinase pathways). This expanded set became the definitive gene list
	for variant retrieval and the subsequent database construction steps.
	
	\subsubsection{Step 3: Variant retrieval from Ensembl BioMart (dbSNP; hsapiens\_snp dataset)}
	After expansion, the STRING output was converted into \textbf{Ensembl Gene Stable IDs} to ensure compatibility with
	Ensembl variant resources. We then queried Ensembl BioMart using the \textbf{hsapiens\_snp} dataset, restricting
	variants to those sourced from \textbf{dbSNP} using:
	\[
	\texttt{variation\_source = dbSNP}.
	\]
	For each gene, we retrieved variant attributes that are directly usable for downstream joins and reporting:
	\textit{Ensembl Gene Stable ID}, gene name, \textit{rsID} (dbSNP identifier), chromosome name, start/end positions,
	alleles, and UCSC identifier (when available).
	
	\begin{figure}[H]
		\centering
		\includegraphics[width=0.92\linewidth]{biomart_variants.png}
		\caption{BioMart attribute selection used in Report 1 for variant retrieval (dbSNP-filtered, hsapiens\_snp).}
		\label{fig:biomart_variants}
	\end{figure}
	
	\paragraph{Finding (Step 3).}
	Direct web-based BioMart queries were not reliable for large gene lists (timeouts, truncated responses, and session
	termination), motivating an automated script-based approach for stable large-scale retrieval.
	
	\subsubsection{Step 4 (automation result): Robust script-based variant acquisition at scale}
	To make variant retrieval reproducible and resilient, Report 1 implemented an automated Python workflow that:
	(i) chunks the gene list, (ii) rotates across multiple Ensembl BioMart mirrors with retries, (iii) checkpoints results
	per chunk, then (iv) merges chunk outputs into a single TSV while preserving exactly one header.
	In the final successful run, \textbf{CHUNK\_SIZE=4} was required for stability, and the full download took
	\textbf{approximately 20 hours} due to BioMart throttling/rate limits and intermittent failures.
	
	\paragraph{Key output artifacts (Step 4).}
	The automated workflow produced:
	\begin{itemize}
		\item \texttt{biomart\_variants.tsv} (merged final variant table),
		\item \texttt{chunks/biomart\_variants\_chunk\_XXXX.tsv} (checkpointed partial results),
		\item \texttt{missing\_gene\_ids.txt} (genes with no returned variant rows),
		\item \texttt{run.log} (execution trace and failure diagnostics).
	\end{itemize}
	
	\paragraph{Why this matters for Report 2.}
	This component is central to the physical design and testing tasks in Report 2 because it defines a
	high-volume, join-heavy table (variants) whose indexing, constraints, and performance properties will strongly affect
	query correctness and runtime in later deliverables.
	
	\subsubsection{Part 3 of Report 1: Transcript-centric genomic coordinates (UCSC Table Browser; hg38.ncbiRefSeq)}
	Report 1 retrieved transcript-level genomic annotations using UCSC Table Browser because the RefSeq track is
	transcript-centric and provides detailed exon structure.
	All UCSC queries were done on \textbf{hg38 (GRCh38)} with:
	\textbf{Group: Genes and Gene Predictions}, \textbf{Track: NCBI RefSeq}, \textbf{Table: ncbiRefSeq}.
	We used gene symbols (mapped from the Ensembl-based gene set) as identifiers and extracted fields necessary for exon
	and coordinate modeling:
	\texttt{name} (RefSeq transcript accession), \texttt{chrom}, \texttt{strand}, \texttt{txStart}, \texttt{txEnd},
	\texttt{exonCount}, \texttt{exonStarts}, \texttt{exonEnds}, and \texttt{name2} (gene symbol).
	Importantly, instead of selecting a single transcript per gene, we retained \textbf{all returned transcripts} to avoid
	arbitrary transcript selection and to preserve alternative splicing information.
	
	\begin{figure}[H]
		\centering
		\includegraphics[width=0.92\linewidth]{Final_Part7.png}
		\caption{UCSC Table Browser configuration snapshot (hg38, RefSeq track, ncbiRefSeq table) used in Report 1.}
		\label{fig:ucsc_config}
	\end{figure}
	
	\paragraph{Finding (UCSC transcript layer).}
	The transcript table provides the ``genomic backbone'' for the database: it enables joins between gene-level entities
	and variant coordinates while supporting transcript-aware questions (e.g., exon count constraints, exon boundary joins,
	and transcript-specific structural reporting).
	
	\subsubsection{Part 4 of Report 1: Gene-to-protein mapping using UniProt}
	To connect genetic identifiers to functional protein entities, Report 1 used UniProt's ID Mapping tool to map the
	phenotype-related gene identifiers into \textbf{UniProt/SwissProt IDs}. The resulting mapping was saved (e.g.,
	\texttt{idmapping.tsv}) and designed to support downstream analyses that require protein attributes such as protein
	length, protein function descriptions, subcellular location, and structural references (e.g., PDB IDs).
	
	\subsubsection{Part 5 of Report 1: GEO dataset discovery and GEO2R groundwork}
	Report 1 identified an expression dataset relevant to schwannoma by constructing an advanced search strategy in NCBI
	GEO using:
	\textbf{keyword: Schwannoma}, \textbf{organism: \textit{Homo sapiens}}, and a publication-date window
	(\textit{Jan 1, 2005 to present}). From the candidate datasets, \textbf{GSE39645} was selected due to biological
	relevance and compatibility with GEO2R-based differential expression analysis.
	This established the foundation for later steps that require integrating log fold-changes and adjusted p-values into
	the relational schema and performing phenotype-relevant filtering and visualization.
	
	\subsection{What Report 1 delivered to enable Report 2}
	Report 1 produced the complete end-to-end blueprint and the core data files needed to proceed with physical design,
	testing, and visualization:
	\begin{enumerate}
		\item A phenotype-driven gene list from BioMart (stable identifiers + coordinates).
		\item A STRING-expanded PPI network (72 nodes) used to define the project gene universe.
		\item A high-volume dbSNP variant dataset retrieved reproducibly via a robust, mirror-aware, chunked script.
		\item Transcript-level RefSeq annotations (hg38.ncbiRefSeq) including exon structure.
		\item UniProt mappings connecting genes to protein entities (supporting later protein-attribute joins).
		\item A selected GEO dataset (GSE39645) and the plan to integrate GEO2R-derived statistics into the schema.
	\end{enumerate}
	
	\paragraph{Summary of key findings.}
	The main scientific/technical takeaways from Report 1 were:
	(1) phenotype-first gene selection provides a reproducible starting point;
	(2) STRING-based expansion adds biologically connected genes that better reflect pathway-level context;
	(3) large-scale variant retrieval from BioMart requires robust automation (chunking, retries, mirror rotation, and
	checkpointing), with \textbf{CHUNK\_SIZE=4} as a stable operating point for this project;
	and (4) transcript-level modeling from UCSC avoids losing splice-isoform structure and better supports exon-aware
	queries required later in the project.
	
	% End of Part 1 recap
	
	\newpage
	
	\section*{Part 2: Physical Database Design and Construction}
	% (Leave blank for now)
	
	\newpage
	
	\section*{Part 3: Database Constraints}
	% (Leave blank for now)
	
	\newpage
	
	\section*{Part 4: Multi-Table SQL Queries}
	% (Leave blank for now)
	
	\newpage
	
	\section*{Part 5: Integrated Query for Variations, Proteins, and Expression}
	% (Leave blank for now)
	
	\newpage
	
	\section*{Part 6: Query for Genes with PDB IDs and Exon Constraints}
	% (Leave blank for now)
	
	\newpage
	
	\section*{Part 7: GEO2R Analysis and Downstream Queries}
	% (Leave blank for now)
	
	\newpage
	
	\section*{Part 8: Protein Interaction Network Visualization}
	% (Leave blank for now)
	
\end{document}
